\section{Method}
\label{sec:method}

\subsection{Looping with a \texorpdfstring{$\beta$}{beta}-residual}

One block $f$ is applied $k$ times with a convex residual,
\begin{equation}
  h^{(i+1)} \;=\; \beta\, f\bigl(h^{(i)}\bigr) + (1-\beta)\, h^{(i)},
  \qquad i = 0,\dots,k-1 .
  \label{eq:loop}
\end{equation}
At $k=1,\beta=1$ this is exactly the unlooped model, which gives us an exact
identity check: a $k=1$, $\beta=1$ looped build must reproduce the baseline
bit-for-bit, and we test it (\texttt{test\_loop\_identity.py}).

\subsection{Choosing \texorpdfstring{$\beta$}{beta}: the Euler reading}
\label{sec:method:euler}

Equation~\eqref{eq:loop} is one explicit Euler step of size $\beta$ on the ODE
\begin{equation}
  \dot h \;=\; f(h) - h .
  \label{eq:ode}
\end{equation}
Under this reading $k$ iterations integrate \eqref{eq:ode} over a total
\emph{flow time}
\[
  T \;=\; k\beta ,
\]
and the choice of $\beta$ becomes a question about what one wants to hold
fixed as $k$ grows. Setting $\beta = 1/k$ keeps $T=1$: deeper models take
\emph{finer} steps along the same trajectory, so $k$ trades discretisation
error for compute and nothing else changes. Setting $\beta = 1/\sqrt{k}$ gives
$T = \sqrt{k}$, which overshoots --- the deeper model integrates further, not
more accurately.

This is not a post-hoc rationalisation: we pre-registered $\beta=1/k$ as the
principled schedule and $\beta=1/\sqrt{k}$ as the control, and the control
loses (\texttt{R1B05}, $-0.39\,\%$, Table~\ref{tab:arms}). The flow-time
reading is the only part of the looping story that our experiments support.

\subsection{Three non-architectural controls}
\label{sec:method:levers}

The core of the experimental design. Each lever changes the data or the
recurrence, adds no depth, and is measured against \texttt{D1}, a baseline
trained on the \emph{same} corpus so that exactly one variable differs.

\paragraph{Window augmentation (\texttt{--augment}).} Random transformations
applied to the training window before normalisation. A regularisation lever:
it should help most where the training corpus is least representative, and its
effect should not depend on horizon.

\paragraph{Negative eigenvalues (\texttt{--allow-neg-eigval}).} The delta-rule
gate is doubled, $\beta \leftarrow 2\beta$, so $\beta \in (0,2)$ and the
transition eigenvalue $1-\beta$ may be negative --- a reflection along the key
direction rather than a contraction (\S\ref{sec:background}). This is a change
to the \emph{reachable state transitions}, not to the optimisation, so we
pre-registered a specific prediction that distinguishes it from the other two
levers: \emph{its advantage should grow with horizon}, because reflections are
what let a recurrence carry alternating or periodic structure forward over many
steps. \S\ref{sec:results} reports that this prediction is the only one that
came true.

\paragraph{Exact-once epochs (\texttt{--sampler enumerated}).} The original
sampler draws windows with replacement, so within one nominal epoch some
windows are seen several times and others never. The v2 sampler replaces this
with a Feistel permutation with cycle-walking, which enumerates each window
exactly once per epoch without storing the permutation, interleaved with the
synthetic corpus by a Bresenham rule that makes the real/synthetic count exact
rather than merely correct in expectation. Because the underlying PRNG is
counter-based (Philox), sample $n$ is a pure function of $(n, \text{seed})$,
which makes the stream independent of the worker count --- the property that
lets us verify two runs saw identical data by comparing a single
\texttt{stream\_sha256}.

\subsection{A quantile head}
\label{sec:method:quantile}

MASE scores the median path only. To measure calibration we attach a
$9$-quantile head to the frozen backbone's output projection and train it with
the pinball loss. Two implementation facts decide whether the numbers mean
anything. First, \texttt{out\_proj} is not the last operation: the decoder
re-applies the window's min--max scale afterwards, so a head trained on
\texttt{out\_proj} outputs lives in normalised space and must be denormalised
per window before the loss. Second, at $q=0.5$ the pinball loss equals exactly
half the MAE, which gives a \emph{step-0 control}: before any training, the
head initialised from the point forecast must reproduce half the MAE to within
a fraction of a percent. Our first implementation failed this control by seven
orders of magnitude and trained for $4000$ steps on a meaningless objective;
the control is now a hard exit rather than a warning.
